\documentclass[12pt,a4paper]{article}

\usepackage[utf8]{inputenc}
\usepackage[T1]{fontenc}
\usepackage{geometry}
\usepackage{hyperref}
\usepackage{enumitem}
\usepackage{xcolor}
\usepackage{listings}
\usepackage{longtable}

\geometry{margin=2.5cm}

\lstdefinestyle{code}{
    basicstyle=\ttfamily\small,
    breaklines=true,
    frame=single,
    columns=fullflexible
}

\title{SRS / SDD \\ T17 - Analiza Semantica Cod Sursa}
\author{Proiect PCD}
\date{2026}

\begin{document}

\maketitle
\tableofcontents
\newpage

\section{Introducere}

Acest document contine specificatia cerintelor software si descrierea de proiectare pentru aplicatia T17 - Analiza Semantica Cod Sursa. Proiectul este realizat pentru disciplina Programare Concurenta si Distribuita si are ca scop implementarea unei aplicatii client-server care permite incarcarea unor fisiere sursa C, analiza lor cu ajutorul bibliotecii libclang si generarea unor rapoarte text.

Aplicatia foloseste comunicare prin socket-uri INET si UNIX, procese multiple, autentificare, transfer de fisiere, generare de rapoarte, operatii de administrare, monitorizare cu inotify, REST API minimal si SOAP API minimal.

\section{Scopul proiectului}

Scopul proiectului este realizarea unui sistem distribuit local care permite:

\begin{itemize}
    \item incarcarea unui fisier sursa C de la client catre server;
    \item analiza fisierului folosind libclang;
    \item generarea unui raport text;
    \item descarcarea raportului de catre client;
    \item administrarea serverului printr-un client separat;
    \item folosirea socket-urilor INET si UNIX;
    \item folosirea mecanismelor \texttt{select()}, \texttt{fork()}, \texttt{exec()}, \texttt{pipe()} si \texttt{inotify};
    \item expunerea unor date prin REST API;
    \item expunerea unor date prin SOAP API si WSDL.
\end{itemize}

\section{Domeniul aplicatiei}

Aplicatia este destinata analizei statice a fisierelor C. Sistemul nu executa codul primit de la client, ci il parseaza si extrage informatii relevante despre structura si posibile probleme din cod.

Analiza include:

\begin{itemize}
    \item numar de functii;
    \item numar de variabile;
    \item numar de instructiuni \texttt{if};
    \item numar de instructiuni \texttt{for};
    \item numar de instructiuni \texttt{while};
    \item diagnostice generate de libclang;
    \item mesaje de tip warning sau error.
\end{itemize}

\section{Descriere generala}

Sistemul are o arhitectura client-server. Clientii normali comunica cu serverul printr-un socket INET TCP pe adresa \texttt{127.0.0.1:8080}. Clientul admin comunica separat printr-un socket UNIX local, folosind calea \texttt{/tmp/t17\_admin.sock}.

Serverul asteapta conexiuni folosind \texttt{select()} pentru a monitoriza simultan socket-ul TCP si socket-ul UNIX. Pentru fiecare conexiune acceptata, serverul creeaza un proces copil cu \texttt{fork()}, iar procesul copil trateaza clientul respectiv.

Pe langa comunicarea prin socket-uri, proiectul include:

\begin{itemize}
    \item REST API minimal pe \texttt{127.0.0.1:8081};
    \item SOAP API minimal pe \texttt{127.0.0.1:8082};
    \item WSDL disponibil la \texttt{http://127.0.0.1:8082/wsdl}.
\end{itemize}

Aceste API-uri nu inlocuiesc serverul principal, ci expun prin HTTP datele generate de server: status, statistici, rapoarte, upload-uri, joburi, useri si loguri.

\section{Cerinte functionale}

\subsection{Cerinte pentru analiza codului}

\begin{enumerate}[label=CF\arabic*.]
    \item Aplicatia trebuie sa permita selectarea unui fisier sursa C pentru analiza.
    \item Aplicatia trebuie sa foloseasca libclang pentru parsarea codului sursa.
    \item Aplicatia trebuie sa extraga statistici despre functii, variabile si structuri de control.
    \item Aplicatia trebuie sa afiseze diagnosticele generate de libclang.
    \item Aplicatia trebuie sa genereze un raport text pentru fiecare fisier analizat.
    \item Aplicatia trebuie sa salveze rapoartele in directorul \texttt{reports/}.
\end{enumerate}

\subsection{Cerinte pentru server}

\begin{enumerate}[label=CF\arabic*., resume]
    \item Serverul trebuie sa porneasca pe socket INET TCP la adresa \texttt{127.0.0.1:8080}.
    \item Serverul trebuie sa creeze un socket UNIX pentru administrare la \texttt{/tmp/t17\_admin.sock}.
    \item Serverul trebuie sa foloseasca \texttt{select()} pentru a astepta conexiuni pe ambele socket-uri.
    \item Serverul trebuie sa foloseasca \texttt{fork()} pentru tratarea concurenta a clientilor.
    \item Serverul trebuie sa autentifice utilizatorii folosind fisierul \texttt{config/users.cfg}.
    \item Serverul trebuie sa accepte comenzi de upload si download pentru clientii normali.
    \item Serverul trebuie sa execute analyzer-ul folosind \texttt{fork()}, \texttt{exec()} si \texttt{pipe()}.
    \item Serverul trebuie sa scrie loguri in \texttt{logs/server.log}.
    \item Serverul trebuie sa mentina statistici in \texttt{logs/stats.txt}.
    \item Serverul trebuie sa mentina istoricul joburilor in \texttt{logs/jobs.log}.
\end{enumerate}

\subsection{Cerinte pentru clientul normal}

\begin{enumerate}[label=CF\arabic*., resume]
    \item Clientul normal trebuie sa se conecteze prin socket INET TCP.
    \item Clientul normal trebuie sa trimita automat o comanda de autentificare.
    \item Clientul normal trebuie sa poata incarca fisiere C catre server.
    \item Clientul normal trebuie sa primeasca rezultatul analizei de la server.
    \item Clientul normal trebuie sa poata descarca rapoarte generate anterior.
\end{enumerate}

\subsection{Cerinte pentru clientul admin}

\begin{enumerate}[label=CF\arabic*., resume]
    \item Clientul admin trebuie sa se conecteze prin socket UNIX.
    \item Clientul admin trebuie sa foloseasca rolul \texttt{admin}.
    \item Clientul admin trebuie sa poata vedea statisticile serverului.
    \item Clientul admin trebuie sa poata vedea logurile serverului.
    \item Clientul admin trebuie sa poata lista fisierele incarcate.
    \item Clientul admin trebuie sa poata lista rapoartele generate.
    \item Clientul admin trebuie sa poata vedea statusul serverului.
    \item Clientul admin trebuie sa poata sterge logurile.
    \item Clientul admin trebuie sa poata sterge un raport.
    \item Clientul admin trebuie sa poata lista utilizatorii configurati.
    \item Clientul admin trebuie sa poata vedea joburile de analiza.
\end{enumerate}

\subsection{Cerinte pentru monitorizare}

\begin{enumerate}[label=CF\arabic*., resume]
    \item Aplicatia trebuie sa includa un program separat care foloseste \texttt{inotify}.
    \item Programul de monitorizare trebuie sa urmareasca directorul \texttt{reports/}.
    \item Programul trebuie sa afiseze evenimente precum creare, modificare, scriere si stergere de rapoarte.
\end{enumerate}

\subsection{Cerinte pentru REST API}

\begin{enumerate}[label=CF\arabic*., resume]
    \item Aplicatia trebuie sa includa un REST API minimal.
    \item REST API-ul trebuie sa ruleze local pe \texttt{127.0.0.1:8081}.
    \item REST API-ul trebuie sa expuna endpoint-uri pentru health, statistici, rapoarte, upload-uri, joburi, useri si loguri.
    \item REST API-ul trebuie sa includa o interfata web minimala la \texttt{/ui}.
    \item REST API-ul trebuie sa aiba specificatie OpenAPI in \texttt{docs/openapi.yaml}.
\end{enumerate}

\subsection{Cerinte pentru SOAP API}

\begin{enumerate}[label=CF\arabic*., resume]
    \item Aplicatia trebuie sa includa un SOAP API minimal.
    \item SOAP API-ul trebuie sa ruleze local pe \texttt{127.0.0.1:8082}.
    \item SOAP API-ul trebuie sa expuna endpoint-ul \texttt{/soap}.
    \item SOAP API-ul trebuie sa ofere un document WSDL la \texttt{/wsdl}.
    \item SOAP API-ul trebuie sa suporte operatii pentru health, statistici, rapoarte, upload-uri, joburi, useri si loguri.
\end{enumerate}

\section{Cerinte nefunctionale}

\begin{enumerate}[label=CNF\arabic*.]
    \item Codul principal trebuie sa fie scris in C.
    \item Proiectul trebuie sa compileze fara warning-uri folosind \texttt{-Wall -Wextra -Werror -pedantic}.
    \item Comunicarea client-server trebuie sa fie stabila pentru fisiere de dimensiuni moderate.
    \item Serverul trebuie sa poata trata mai multi clienti prin procese copil.
    \item Rapoartele si logurile trebuie salvate local in directoare separate.
    \item Clientul admin trebuie separat logic de clientul normal.
    \item Documentatia oficiala trebuie scrisa in LaTeX.
    \item REST API-ul trebuie sa poata fi testat prin \texttt{scripts/test\_rest.sh}.
    \item SOAP API-ul trebuie sa poata fi testat prin \texttt{scripts/test\_soap.sh}.
\end{enumerate}

\section{Arhitectura sistemului}

Arhitectura generala este urmatoarea:

\begin{lstlisting}[style=code]
Client C normal      \
                     -> INET TCP 127.0.0.1:8080 -> Server -> Analyzer -> Report
Client Python        /

Admin Client C       -> UNIX socket /tmp/t17_admin.sock -> Server

Watch Reports        -> inotify -> reports/

REST API             -> HTTP 127.0.0.1:8081 -> logs/, reports/, uploads/, config/

SOAP API             -> HTTP 127.0.0.1:8082 -> logs/, reports/, uploads/, config/
\end{lstlisting}

Serverul principal are doua puncte de intrare:

\begin{itemize}
    \item socket INET TCP pentru clientii normali;
    \item socket UNIX local pentru clientul admin.
\end{itemize}

REST API-ul si SOAP API-ul sunt componente separate, folosite pentru a expune prin HTTP informatii deja generate de serverul principal.

\section{Componente software}

\subsection{\texttt{src/analyzer.c}}

Componenta care foloseste libclang pentru analiza fisierelor C. Primeste fisierul prin argumente de linie de comanda si afiseaza rezultatul la \texttt{stdout}. Serverul captureaza acest output prin pipe.

\subsection{\texttt{src/server.c}}

Serverul principal. Creeaza socket-ul INET, socket-ul UNIX, foloseste \texttt{select()}, accepta clienti, foloseste \texttt{fork()}, autentifica utilizatori, primeste fisiere, ruleaza analyzer-ul si gestioneaza comenzile admin.

\subsection{\texttt{src/client.c}}

Clientul C normal pentru upload de fisiere si download de rapoarte.

\subsection{\texttt{src/admin\_client.c}}

Clientul admin C, conectat prin socket UNIX. Permite operatii precum statistici, loguri, listare rapoarte, listare useri, joburi si stergere rapoarte.

\subsection{\texttt{python\_client/client.py}}

Client alternativ in Python pentru upload si download prin acelasi protocol TCP.

\subsection{\texttt{src/watch\_reports.c}}

Program separat care foloseste \texttt{inotify} pentru monitorizarea directorului \texttt{reports/}.

\subsection{\texttt{rest\_api/app.py}}

REST API minimal implementat in Python. Ruleaza pe \texttt{127.0.0.1:8081} si expune date in format JSON.

\subsection{\texttt{soap\_api/soap\_server.py}}

SOAP API minimal implementat in Python. Ruleaza pe \texttt{127.0.0.1:8082}, expune endpoint-ul \texttt{/soap} si documentul WSDL la \texttt{/wsdl}.

\subsection{\texttt{scripts/test\_rest.sh}}

Script pentru testarea endpoint-urilor REST cu \texttt{curl}.

\subsection{\texttt{scripts/test\_soap.sh}}

Script pentru testarea operatiilor SOAP cu request-uri XML trimise prin \texttt{curl}.

\section{Structura proiectului}

\begin{lstlisting}[style=code]
Proiect_PCD/
|-- Makefile
|-- README.md
|-- README.tex
|-- config/
|   |-- app.cfg
|   |-- users.cfg
|-- docs/
|   |-- SRS_SDD.tex
|   |-- protocol.tex
|   |-- openapi.yaml
|-- python_client/
|   |-- client.py
|-- rest_api/
|   |-- app.py
|-- soap_api/
|   |-- soap_server.py
|-- scripts/
|   |-- test_rest.sh
|   |-- test_soap.sh
|-- src/
|   |-- analyzer.c
|   |-- main.c
|   |-- server.c
|   |-- client.c
|   |-- admin_client.c
|   |-- watch_reports.c
|-- tests/
|   |-- sample.c
|   |-- bad.c
|   |-- bad_syntax.c
|   |-- bad_type.c
|   |-- bad_uninitialized.c
|   |-- good_math.c
|   |-- good_nested.c
|   |-- good_while.c
\end{lstlisting}

\section{Protocol socket}

Protocolul socket este text-based. Fiecare client trimite mai intai:

\begin{lstlisting}[style=code]
LOGIN <username> <password>
\end{lstlisting}

Clientii normali pot folosi:

\begin{lstlisting}[style=code]
UPLOAD <filename> <size>
DOWNLOAD_REPORT <report_name>
QUIT
\end{lstlisting}

Clientul admin poate folosi:

\begin{lstlisting}[style=code]
STATS
LOGS
LIST_UPLOADS
LIST_REPORTS
SERVER_STATUS
CLEAR_LOGS
DELETE_REPORT <report_name>
LIST_USERS
JOBS
QUIT
\end{lstlisting}

\section{REST API}

REST API-ul ruleaza separat de serverul TCP/UNIX:

\begin{lstlisting}[style=code]
python3 rest_api/app.py
\end{lstlisting}

Adresa locala:

\begin{lstlisting}[style=code]
http://127.0.0.1:8081
\end{lstlisting}

Endpoint-uri disponibile:

\begin{itemize}
    \item \texttt{GET /health};
    \item \texttt{GET /stats};
    \item \texttt{GET /reports};
    \item \texttt{GET /reports/<name>};
    \item \texttt{GET /uploads};
    \item \texttt{GET /jobs};
    \item \texttt{GET /users};
    \item \texttt{GET /logs};
    \item \texttt{GET /ui}.
\end{itemize}

REST API-ul intoarce raspunsuri JSON.

\section{SOAP API}

SOAP API-ul ruleaza separat:

\begin{lstlisting}[style=code]
python3 soap_api/soap_server.py
\end{lstlisting}

Endpoint SOAP:

\begin{lstlisting}[style=code]
http://127.0.0.1:8082/soap
\end{lstlisting}

Document WSDL:

\begin{lstlisting}[style=code]
http://127.0.0.1:8082/wsdl
\end{lstlisting}

Operatii SOAP suportate:

\begin{itemize}
    \item \texttt{Health};
    \item \texttt{Stats};
    \item \texttt{Reports};
    \item \texttt{Uploads};
    \item \texttt{Jobs};
    \item \texttt{Users};
    \item \texttt{Logs}.
\end{itemize}

SOAP API-ul intoarce raspunsuri XML in format SOAP Envelope.

\section{Fluxuri de executie}

\subsection{Upload si analiza fisier}

\begin{enumerate}
    \item Clientul se conecteaza la server prin TCP.
    \item Clientul trimite automat comanda \texttt{LOGIN}.
    \item Serverul verifica utilizatorul in \texttt{config/users.cfg}.
    \item Clientul trimite comanda \texttt{UPLOAD}.
    \item Serverul salveaza fisierul in \texttt{uploads/}.
    \item Serverul scrie jobul cu status \texttt{QUEUED}.
    \item Serverul ruleaza analyzer-ul cu \texttt{fork()}, \texttt{exec()} si \texttt{pipe()}.
    \item Serverul scrie status \texttt{RUNNING}, apoi \texttt{DONE} sau \texttt{FAILED}.
    \item Serverul salveaza raportul in \texttt{reports/}.
    \item Serverul trimite rezultatul catre client.
\end{enumerate}

\subsection{Administrare server}

\begin{enumerate}
    \item Admin clientul se conecteaza la \texttt{/tmp/t17\_admin.sock}.
    \item Admin clientul trimite login cu rol de admin.
    \item Serverul verifica rolul.
    \item Adminul trimite una dintre comenzile suportate.
    \item Serverul intoarce rezultatul.
\end{enumerate}

\subsection{Monitorizare rapoarte}

\begin{enumerate}
    \item Utilizatorul porneste \texttt{watch\_reports}.
    \item Programul adauga un watch cu \texttt{inotify\_add\_watch()} pe \texttt{reports/}.
    \item Cand apare un raport nou, programul afiseaza evenimentul.
\end{enumerate}

\subsection{Acces REST}

\begin{enumerate}
    \item Utilizatorul porneste REST API-ul.
    \item REST API-ul ruleaza pe \texttt{127.0.0.1:8081}.
    \item Utilizatorul acceseaza endpoint-uri prin browser sau \texttt{curl}.
    \item API-ul citeste fisierele generate de server si intoarce JSON.
\end{enumerate}

\subsection{Acces SOAP}

\begin{enumerate}
    \item Utilizatorul porneste SOAP API-ul.
    \item SOAP API-ul ruleaza pe \texttt{127.0.0.1:8082}.
    \item Utilizatorul trimite cereri XML SOAP catre \texttt{/soap}.
    \item API-ul citeste fisierele generate de server si intoarce raspuns SOAP XML.
    \item WSDL-ul poate fi accesat la \texttt{/wsdl}.
\end{enumerate}

\section{Biblioteci si apeluri folosite}

\begin{longtable}{|p{5cm}|p{9cm}|}
\hline
Componenta & Utilizare \\
\hline
\texttt{socket()} & creare socket INET si UNIX \\
\hline
\texttt{bind()} & asociere socket cu adresa si port/cale locala \\
\hline
\texttt{listen()} & trecerea socket-ului in modul de asteptare conexiuni \\
\hline
\texttt{accept()} & acceptarea unui client \\
\hline
\texttt{select()} & asteptare simultana pe socket TCP si socket UNIX \\
\hline
\texttt{fork()} & creare proces copil pentru fiecare client \\
\hline
\texttt{exec()} & rularea analyzer-ului \\
\hline
\texttt{pipe()} & capturarea rezultatului analyzer-ului \\
\hline
\texttt{waitpid()} & asteptarea proceselor copil \\
\hline
\texttt{inotify} & monitorizarea directorului de rapoarte \\
\hline
\texttt{libclang} & analiza statica a codului C \\
\hline
\texttt{libconfig} & citirea configuratiei aplicatiei \\
\hline
\texttt{http.server} & REST API si SOAP API minimal in Python \\
\hline
\texttt{curl} & testarea endpoint-urilor REST si SOAP \\
\hline
\end{longtable}

\section{Date si fisiere generate la runtime}

\begin{itemize}
    \item \texttt{uploads/} - fisierele incarcate de clienti;
    \item \texttt{reports/} - rapoartele generate;
    \item \texttt{logs/server.log} - logurile serverului;
    \item \texttt{logs/stats.txt} - statistici despre fisiere analizate;
    \item \texttt{logs/jobs.log} - istoricul joburilor de analiza;
    \item \texttt{/tmp/t17\_admin.sock} - socket UNIX pentru admin.
\end{itemize}

\section{Testare}

\subsection{Compilare}

\begin{lstlisting}[style=code]
make clean
make
\end{lstlisting}

\subsection{Test server si clienti}

\begin{lstlisting}[style=code]
./server
./client tests/sample.c
./client tests/bad_type.c
./client download sample_report.txt
python3 python_client/client.py upload tests/good_math.c
python3 python_client/client.py download good_math_report.txt
./admin_client
\end{lstlisting}

\subsection{Test inotify}

\begin{lstlisting}[style=code]
./watch_reports
./client tests/sample.c
\end{lstlisting}

\subsection{Test REST API}

\begin{lstlisting}[style=code]
python3 rest_api/app.py
./scripts/test_rest.sh
\end{lstlisting}

\subsection{Test SOAP API}

\begin{lstlisting}[style=code]
python3 soap_api/soap_server.py
./scripts/test_soap.sh
curl http://127.0.0.1:8082/wsdl
\end{lstlisting}

\section{Limitari curente}

\begin{itemize}
    \item analiza este concentrata pe fisiere C;
    \item protocolul socket este simplu, text-based;
    \item REST API-ul si SOAP API-ul sunt minimale si citesc date din fisiere locale;
    \item job tracking-ul este salvat in fisier text;
    \item serverul este destinat pentru testare locala;
    \item autentificarea este simpla si foloseste fisier text.
\end{itemize}

\section{Directii viitoare}

\begin{itemize}
    \item REST API mai complex cu metode POST si DELETE;
    \item SOAP API mai complet;
    \item coada asincrona reala pentru joburi;
    \item baza de date pentru useri, rapoarte si joburi;
    \item hashing pentru parole;
    \item suport mai bun pentru C++;
    \item interfata web mai dezvoltata.
\end{itemize}

\section{Concluzie}

Proiectul implementeaza o aplicatie client-server functionala pentru analiza semantica a codului sursa C. Sistemul foloseste socket-uri INET pentru clientii normali, socket UNIX pentru administrare, \texttt{select()} pentru monitorizarea mai multor socket-uri, \texttt{fork()} pentru concurenta, \texttt{exec()} si \texttt{pipe()} pentru rularea analyzer-ului, \texttt{inotify} pentru monitorizarea rapoartelor generate, REST API si SOAP API pentru expunerea datelor prin HTTP.

Prin aceste functionalitati, proiectul acopera cerinte importante din zona de programare concurenta si distribuita.

\end{document}